\subsection{Two compact Wigner chains with a uniform lower bound}
\label{sec:archimedean-wigner-two-chain}

We isolate the compact representation calculation used when the two
compact types coincide. The classical input is the orthonormal
Gelfand--Tsetlin branching theorem for \(U(k)\), together with the
one-step Clebsch--Gordan coefficient formulas for tensoring with the
standard representation and its dual. We use the normalization in
Appendix~A of Grinko--Burchardt--Ozols.\cite{grinko-burchardt-ozols-wigner}
The estimates below are proved from those formulas; no assertion about
a formal proof replay is involved.

Write \(k=n+1\ge2\). A highest weight
\(\lambda_0\ge\cdots\ge\lambda_n\), with every \(\lambda_j\in\mathbb Z\),
has shifted labels \(L_j=\lambda_j-j\).
Thus \(L_0>\cdots>L_n\) are integers. A summand of its
restriction to the subgroup \(U(n)\) fixing the last standard vector
has shifted labels \(M_0,\ldots,M_{n-1}\), with
\[
                 L_{j+1}<M_j\le L_j\quad(0\le j<n).
                                                        \tag{WC1}
\]
The restriction is an orthogonal, multiplicity-free sum of these
summands. In these labels, the squares of the two normalized one-step
coefficients at the last standard vector are
\[
 \begin{split}
 w_+(L,M,l)&=
 \frac{\prod_{j=0}^{n-1}|M_j-L_l-1|}
      {\prod_{i\ne l}|L_i-L_l|},\\
 w_-(L,M,l)&=
 \frac{\prod_{j=0}^{n-1}|M_j-L_l|}
      {\prod_{i\ne l}|L_i-L_l|}.
 \end{split}                                           \tag{WC2}
\]
Here the raising or lowering of the \(l\)-th label is required to
remain dominant. An absent branch has coefficient zero.

\begin{lemma}[No rank loss in a nonzero one-step coefficient]
\label{lem:wigner-one-step-lower}
Assume (WC1), \(D\ge1\),
\(L_0-L_l\le D\), and \(L_l-L_n\le D\).
If the numerator defining \(w_+\), respectively \(w_-\), is nonzero,
then the corresponding squared coefficient \(w_+\) or \(w_-\)
is at least \(D^{-2}\).
\end{lemma}

\begin{proof}
Set \(d_i=\max(1,|L_i-L_l|)\). In particular \(d_l=1\),
and strict integral ordering gives \(d_i=|L_i-L_l|\) for \(i\ne l\).
For either numerator in (WC2), its factor with index \(j\) is at least
\(d_{j+1}\) when \(j<l\), and at least \(d_j\) when \(j\ge l\).
For example, when \(j<l\), integrality gives
\(M_j\ge L_{j+1}+1\). For \(w_+\) this gives
\(M_j-L_l-1\ge L_{j+1}-L_l\); if \(j+1=l\), the
nonzero-numerator assumption improves the possible zero to at least
one. For \(w_-\) the same lower bound is immediate. When \(j\ge l\),
use \(M_j\le L_j\le L_l\); the only possible zero requiring the
same improvement is now in \(w_-\) at \(j=l\).

The product of these lower bounds is
\[
 P_l=\left(\prod_{j<l}d_{j+1}\right)
                   \left(\prod_{j\ge l}d_j\right).
\]
The identity
\[
             \prod_{i\ne l}d_i=d_0d_nP_l
                                                        \tag{WC3}
\]
holds also for \(l=0,n\), with empty products equal to one.
All interior factors cancel and \(d_l=1\). Thus either nonzero
coefficient is at least \((d_0d_n)^{-1}\ge D^{-2}\).
\end{proof}

\begin{theorem}[Two-chain contraction inequality]
\label{thm:wigner-two-chain-lower}
Let \(V_\lambda\) be the irreducible unitary representation just
described. Choose \(0\le i<n\), \(N\in\mathbb Z_{\ge0}\), and
\(b\ge1\) such that
\[
 L_i-L_{i+1}\ge2N+1,\qquad N\le b,
                    \qquad L_0-L_n\le2b .              \tag{WC4}
\]
There are two polynomial maps, of total degree \(N\) in
\(\omega,\overline\omega\),
\[
 \begin{split}
 A_+(\omega)&:V_{\lambda+Ne_{i+1}}\longrightarrow V_\lambda,\\
 A_-(\omega)&:V_{\lambda-Ne_i}\longrightarrow V_\lambda,
 \qquad \omega\in\mathbb C^k,
 \end{split}
\]
equivariant for \(U(k)\) and homogeneous of degree \(N\) under
real scalar multiplication. They satisfy
\[
                  \|A_\pm(\omega)\|_{\rm op}
                                  \le\|\omega\|^N       \tag{WC5}
\]
and, for every \(v\in V_\lambda\),
\[
 \|\omega\|^{2N}\|v\|^2
 \le(3b)^{2N}
       \bigl(\|A_+(\omega)^*v\|^2+
                         \|A_-(\omega)^*v\|^2\bigr).    \tag{WC6}
\]
Neither bound contains an additional factor depending on \(k\) or
the dimensions of the representation spaces.
\end{theorem}

\begin{proof}
At each raising step use the isometric irreducible inclusion
\(V_{\lambda+e_l}\hookrightarrow V_\lambda\otimes\mathbb C^k\)
and contract the standard factor against \(\omega\). At each
lowering step use the corresponding isometric inclusion into
\(V_\lambda\otimes(\mathbb C^k)^*\) and contract against the
dual vector determined by \(\omega\). These contractions are
conjugate linear or linear in \(\omega\), respectively; both are
polynomial in \(\omega,\overline\omega\), equivariant, and of
operator norm at most \(\|\omega\|\).
Compose \(N\) raising steps at index \(i+1\), or \(N\) lowering
steps at index \(i\), in the order from the final shifted space
back to \(V_\lambda\). Condition (WC4) ensures dominance of all
intermediate weights. This constructs the maps and proves (WC5).
For \(N=0\), use the identity twice; (WC6) is then immediate.
We henceforth assume \(N>0\).

First take \(\omega=e_k\). Every step intertwines the subgroup
\(U(k-1)\) fixing \(e_k\); its adjoint therefore preserves the
branch label \(M\). On any branch that survives all stages,
the eigenvalue of \(A_+(e_k)A_+(e_k)^*\), respectively
\(A_-(e_k)A_-(e_k)^*\), is
\[
 \prod_{t=0}^{N-1}w_+(L+te_{i+1},M,i+1),\qquad
 \prod_{t=0}^{N-1}w_-(L-te_i,M,i).                       \tag{WC7}
\]
This follows inductively from (WC2): the adjoint carries the
branch into the branch with the same \(M\), and the next positive
operator is scalar on that branch. If a branch disappears, that
chain kills it at the first zero coefficient. We use only a chain
for which all branches survive.

Fix a branch \(M\) of the original representation. If
\[
                        M_i\ge L_{i+1}+N+1,             \tag{WC8}
\]
then (WC1) remains true after replacing \(L_{i+1}\) by
\(L_{i+1}+t\), for every \(0\le t\le N\).
None of the raising numerators vanishes for \(t<N\).
Indeed the index \(i\) is separated from
\(L_{i+1}+t+1\) by (WC8). For \(j<i\), interlacing gives
\(M_j>L_i\), and for \(j>i\) it gives \(M_j\le L_{i+1}\).
These exclude every other possible zero.

If (WC8) fails, integrality gives
\(M_i\le L_{i+1}+N\). The gap in (WC4) then implies
\[
                       M_i<L_i-N .                     \tag{WC9}
\]
Consequently (WC1) survives every lowering step
\(L_i\mapsto L_i-t\), \(0\le t\le N\), and no lowering
numerator vanishes for \(t<N\). At index \(i\) use (WC9);
at smaller indices use \(M_j>L_i\); at larger indices use
\(M_j\le L_{i+1}\). This proves that at least one of the
two complete chains survives with nonzero coefficients.

At any such stage the changed label differs from its original
value by at most \(N\le b\). Its distance to either endpoint
of the current ordered list is at most \(3b\), since the
original span is at most \(2b\). Lemma~\ref{lem:wigner-one-step-lower}
therefore bounds every factor in the successful product (WC7)
below by \((3b)^{-2}\). On every original branch, the sum of the
two positive operators has eigenvalue at least \((3b)^{-2N}\).
Orthogonality of the branch decomposition proves (WC6) at \(e_k\).

Unitary covariance proves it at every unit vector, since \(U(k)\)
acts transitively on the unit sphere. Positive real homogeneity
proves it at every nonzero vector. At zero both sides vanish
when \(N>0\), completing the proof.
\end{proof}

\paragraph{Relation to the local compact-type convention.}
The compact parameter in the local problem is expressed in coordinates
\(u_{j+1}=L_j+c\), with a common real offset \(c\).
Its differences and its updates by \(\pm N\) are exactly those above;
in particular \(|u_j|\le b\) gives the span bound in (WC4).
This is a change of parameter coordinates, not the assertion that
every real offset defines a character of \(U(k)\). Honest common
character twists are \(\det^q\), \(q\in\mathbb Z\); they may be
applied to both sides of each intertwiner and leave the norms and
coefficient formulas unchanged.
An isometric identification of the positive compact quotient with
\(\mathbb C^k\) replaces \(\|\omega\|^2\) in (WC5)--(WC6) by
its original positive Hermitian form. Thus the estimate is
independent of the choice of orthonormal quotient frame.

This proves the compact two-chain lower bound and the contraction
upper bound used in the equal-type argument. To obtain the full
archimedean estimate one must still transport these maps to the
actual ambient polynomial and radial operators, retain their
normalizations, and justify the uniform differential-word estimate
and removal of the Gram cutoff. The compact calculation alone
does not prove (AP6).
